% Generated draft: requires mathematical and source review.
\documentclass[11pt]{article}
\usepackage[margin=25mm]{geometry}
\usepackage{fontspec}
\setmainfont{Latin Modern Roman}
\usepackage{amsmath,amssymb,mathtools,unicode-math}
\setmathfont{Latin Modern Math}
\usepackage{enumitem,needspace,xurl,hyperref}
\setlength{\parindent}{0pt}
\setlength{\parskip}{5pt}
\setlength{\emergencystretch}{3em}
\newcommand{\N}{\mathbb{N}}
\newcommand{\Z}{\mathbb{Z}}
\newcommand{\Q}{\mathbb{Q}}
\newcommand{\R}{\mathbb{R}}
\newcommand{\C}{\mathbb{C}}
\newcommand{\F}{\mathbb{F}}
\newcommand{\A}{\mathbb{A}}
\newcommand{\PP}{\mathbb{P}}
\newcommand{\E}{\mathbb{E}}
\newcommand{\eps}{\varepsilon}
\newcommand{\dd}{\,\mathrm{d}}
\begin{document}
\section*{UnsolvedMath ID: 8800061}
\textbf{Source problem code:} catalogue entry
\subsection{Definitions and mathematical statement}
In a common prime-order pairing setting, let \(G_1,G_2\) be additive cyclic groups of order \(r\), let \(G_T\) be a multiplicative cyclic group of order \(r\), and let an efficiently computable nondegenerate bilinear map satisfy
\[
e:G_1\times G_2\to G_T,\qquad
 e(aP,bQ)=e(P,Q)^{ab}.
\]
The fixed-first-argument inversion problem is: given \(P\in G_1\setminus\{0\}\) and \(z\in G_T\), find \(Q\in G_2\) such that
\[
e(P,Q)=z.
\]
Nondegeneracy and prime order make this preimage unique. Fixing the second argument gives another version; asking for both arguments to be found is a different formulation.

Lange asks to determine the computational hardness of the intended Pairing Inversion Problem. A precise claim must specify the pairing family and group encodings, fixed/free arguments, instance distribution, size parameter, deterministic/randomized or quantum adversaries, and whether hardness means a reduction to a standard problem, a model lower bound, or an average-case assumption. Nondegeneracy alone is an algebraic property, not a uniform computational hardness statement.

For orientation, if \(Q_0\) generates \(G_2\), a discrete-logarithm solver in \(G_T\) can write \(z=e(P,Q_0)^a\) and return \(aQ_0\). That direction does not show that pairing inversion is as hard as discrete logarithms. A claimed hardness reduction must have its actual direction and hypotheses stated.

The supplied research summary attributes a systematic 2008 study and special-case algorithms to Galbraith--Hess--Vercauteren. It uses broad hardness language while also saying no general polynomial-time algorithm or hardness proof is known. Without a fixed variant, model, and theorem those phrases cannot be reconciled as an unconditional general hardness result; they are retained as a qualified partial-progress assessment.
\subsection{Short English statement}
Determine the complexity of recovering a pairing argument from a target value, with the pairing family, inversion variant, encoding, and hardness notion explicit. Bilinearity and a discrete-log upper reduction alone do not establish a hardness theorem.
\subsection{Sources}
\begin{itemize}
\item Local UnsolvedMath registry, permanent ID \(8800061\), exact source code \texttt{catalogue entry}, ``Hardness of the Pairing Inversion Problem.'' Proposer Tanja Lange.
\item IPAM \emph{Workshop on Number Theory and Cryptography --- Open Problems} (2006), attachment 6239, recorded PDF page 50, printed slide 48; \url{https://www.ipam.ucla.edu/programs/workshops/workshop-i-number-theory-and-cryptography-open-problems/?tab=schedule}. Recorded access 2026-07-29; pdf-and-legacy-ppt visual-audit extraction.
\item Supplied PARTIAL-PROGRESS research summary cites Galbraith--Hess--Vercauteren (2008), without title, publication details, exact variant, or the unnamed factor governing its reported special-case algorithms.
\end{itemize}
\end{document}
